% ============================================================
%  1. PRESENTACIÓN DE LA EMPRESA  (Subdocumento 1 — Formulario T-6)
%  Informe 1 — Presentación de la empresa
% ============================================================
\chapter{Presentación de la Empresa}
\label{sec:empresa}

\section{Identificación y Perfil Corporativo}

\textbf{LafroX} es una firma de ingeniería de software y consultoría tecnológica avanzada
con más de una década de trayectoria en el mercado latinoamericano. Nos especializamos en
el diseño, arquitectura y despliegue de soluciones de software resilientes para
ecosistemas de misión crítica, enfocando nuestro núcleo de \emph{expertise} en la
modernización de la cadena de suministro, la logística de última milla y la distribución
de consumo masivo (FMCG).

Nuestro factor diferenciador corporativo radica en el dominio profundo de arquitecturas
distribuidas (Cloud-Edge Computing), el desarrollo de aplicaciones móviles robustas bajo
la filosofía \emph{offline-first} para entornos hostiles sin cobertura de red, y la
integración de telemetría IoT de alta precisión para la trazabilidad industrial y de
cadena de frío. En LafroX no construimos sistemas genéricos; diseñamos plataformas de
ingeniería robustas, adaptables y tolerantes a fallos, capaces de sostener operaciones
transaccionales ininterrumpidas bajo las condiciones logísticas y geográficas más
extremas del territorio nacional.

\section{Filosofía Estratégica: Misión, Visión y Sostenibilidad}

\subsection{Misión}

Empoderar a las organizaciones logísticas, de transporte y distribución masiva mediante
el diseño de software altamente resiliente, arquitecturas tolerantes a fallos y flujos
automáticos descentralizados. Ayudamos a nuestros clientes a optimizar la eficiencia de
sus activos, garantizar la trazabilidad de extremo a extremo de sus operaciones y
reducir el costo de servir (\emph{Cost to Serve}) capturando el dato preciso en el
origen del hecho operacional, garantizando la continuidad del negocio sin importar el
estado de su conectividad.

\subsection{Visión}

Consolidarnos para el año 2030 como el socio tecnológico de referencia en el Cono Sur
para la transformación digital e inteligente de la cadena de suministro
(\emph{supply chain}). Lideramos la transición de la industria hacia modelos logísticos
eficientes impulsados por la adopción de arquitecturas nativas de la nube, computación en
el borde (\emph{Edge Computing}), analítica predictiva de datos y automatización móvil
bajo los más exigentes estándares globales de ingeniería de software.

\subsection{Compromiso de Sostenibilidad y Gobernanza (ESG)}

LafroX concibe la tecnología moderna no solo como un catalizador de rentabilidad
corporativa, sino también como un habilitador directo para el cumplimiento de las
normativas medioambientales, de responsabilidad social y gobernanza (ESG). Nuestras
arquitecturas y soluciones integran estos compromisos de manera nativa mediante tres
pilares operativos de alto impacto:

\begin{itemize}
  \item \textbf{Reducción activa de la huella de carbono:} implementamos algoritmos
        heurísticos avanzados para la optimización y el ruteo dinámico de flotas. Estos
        modelos calculan las rutas óptimas, minimizando los kilómetros recorridos en
        vacío, mejorando los índices de capacidad de carga por vehículo y reduciendo de
        manera directa las emisiones de gases de efecto invernadero (GEI). Todas nuestras
        estrategias de optimización respetan las directrices, incentivos y metas de la
        Estrategia Nacional de Movilidad Sostenible (ENMS) de Chile.
  \item \textbf{Eficiencia energética y \emph{cloud} verde:} priorizamos el despliegue
        de nuestras soluciones en centros de datos de hiperescala (AWS/Azure) con
        certificaciones internacionales de carbono neutralidad y energía 100\,\%
        renovable. El diseño de nuestro software utiliza arquitecturas \emph{serverless}
        optimizadas que demandan recursos solo bajo demanda real, evitando el consumo
        energético ocioso de los centros de datos \emph{on-premise} tradicionales.
  \item \textbf{Despapelización de procesos operacionales:} reemplazamos los flujos
        documentales físicos mediante la habilitación de firmas electrónicas avanzadas,
        guías de despacho digitales autorizadas por el SII (Documentos Tributarios
        Electrónicos, DTE) y sistemas de pruebas de entrega (\emph{Proof of Delivery},
        PoD) móviles.
\end{itemize}

\section{Trayectoria y Casos de Éxito}

La experiencia de LafroX ha sido validada en operaciones logísticas complejas de alta
intensidad a nivel nacional. A continuación se destacan dos hitos tecnológicos
representativos de nuestras capacidades.

\subsection{Modernización de última milla y operación desconectada}

\begin{itemize}
  \item \textbf{Cliente referencial:} Embotelladora Andina (industria de bebidas).
  \item \textbf{Desafío técnico:} la flota de reparto capilar enfrentaba severas pérdidas
        de información en tiempo real, retrasos en la rendición de cuentas y altas tasas
        de reentrega debido a la intermitencia o nulidad de la cobertura móvil en rutas
        rurales y de geografía compleja.
  \item \textbf{Solución implementada:} arquitectura híbrida \emph{Edge-Cloud} acoplada a
        una suite de aplicaciones móviles nativas bajo el patrón \emph{offline-first}.
        La plataforma permite a los conductores descargar de forma segura su ruta de
        despacho diaria cifrada en una base de datos local (encriptación en disco). Toda
        transacción comercial, registro de cobranza, devoluciones y captura de firmas de
        prueba de entrega se ejecutan y validan localmente en el móvil sin requerir
        conectividad.
  \item \textbf{Impacto operativo:} reducción a cero de la pérdida de guías de despacho
        físicas, conciliación y sincronización asíncrona determinista con el ERP en menos
        de 5 minutos al retornar al centro de distribución, y un aumento sostenido del
        8\,\% en el indicador OTIF (\emph{On-Time In-Full}).
\end{itemize}

\subsection{Trazabilidad sanitaria lote-a-lote y cadena de frío IoT}

\begin{itemize}
  \item \textbf{Cliente referencial:} Agrosuper (producción y distribución de alimentos).
  \item \textbf{Desafío técnico:} garantizar la trazabilidad milimétrica de productos
        perecibles desde la recepción en el centro de distribución hasta el canal
        tradicional, asegurando el control inmutable de los rangos de temperatura para
        cumplir con la auditoría sanitaria y mitigar multas o mermas de stock.
  \item \textbf{Solución implementada:} plataforma de trazabilidad basada en el estándar
        internacional GS1, integrada de extremo a extremo con telemetría IoT.
        Implementamos hardware industrial especializado (\emph{rugged}) apto para operar
        de forma continua en cámaras de congelación extrema ($-$22\,°C) y conectamos
        termógrafos vehiculares inalámbricos a un motor de reglas en la nube para alertar
        excursiones de temperatura críticas en tiempo real.
  \item \textbf{Impacto operativo:} mitigación drástica de pérdidas de stock por quiebre
        de cadena de frío y reducción del tiempo de respuesta ante auditorías de
        trazabilidad sanitarias de días hábiles a escasos minutos.
\end{itemize}

\section{Credenciales Organizacionales y Madurez TI}

LafroX certifica que sus procesos metodológicos para el ciclo de vida del desarrollo de
software (SDLC) y la operación de infraestructuras TI se guían por los más exigentes
estándares internacionales:

\begin{itemize}
  \item \textbf{CMMI-DEV Nivel 3 (Definido):} prácticas organizacionales de ingeniería,
        gestión de requisitos, verificación, validación e integración de sistemas de
        software formalizadas y documentadas.
  \item \textbf{ISO 9001:} Sistema de Gestión de Calidad (SGC) certificado, asegurando
        la mejora continua en todos los entregables técnicos.
  \item \textbf{ISO/IEC 20000-1 / ITIL v4:} ciclo de soporte técnico, gestión de
        incidentes y mesa de ayuda alineados a la norma de gestión de servicios de TI y
        a las mejores prácticas de ITIL v4, garantizando el cumplimiento de los Acuerdos
        de Nivel de Servicio (SLA).
  \item \textbf{ISO/IEC 27001:} protección de los activos de información mediante
        controles rigurosos, cifrado robusto para datos en reposo (AES-256) y en tránsito
        (TLS 1.3), políticas Zero-Trust y monitoreo de intrusiones.
  \item \textbf{DevSecOps:} metodologías ágiles escaladas (Scrum y SAFe) con tuberías de
        integración, entrega y seguridad continua; cada pieza de código se somete a
        escaneos automáticos de vulnerabilidades antes de ser desplegada en producción.
\end{itemize}

\section{Capacidad Instalada para Proyectos Complejos}

\begin{itemize}
  \item \textbf{Socio AWS Partner Network (APN):} red global de socios tecnológicos de
        Amazon Web Services, con arquitectos certificados en microservicios
        \emph{serverless}, persistencia políglota distribuida e ingesta masiva de datos
        en tiempo real.
  \item \textbf{Laboratorio de simulación Edge-to-Cloud y pruebas offline:} centro de
        experimentación con cámaras y jaulas de Faraday para simular intermitencia
        extrema y pérdida total de enlaces (GPRS/3G/4G/5G), certificando la consistencia
        del almacenamiento encriptado local (SQLCipher, AES-256) y la robustez de los
        algoritmos de sincronización asíncrona y resolución de conflictos.
  \item \textbf{Plataforma de despliegue elástico y segregación de ambientes:}
        infraestructura CI/CD automatizada sobre AWS Fargate y Amazon ECS, con cuatro
        entornos independientes y securizados mediante WAF: DEV, QA, STG y PROD.
  \item \textbf{Consola NOC y observabilidad activa:} central de monitoreo 24/7 con
        herramientas de APM y observabilidad (Datadog/CloudWatch), respaldando un
        compromiso de RTO de 30 minutos y RPO de 5 minutos ante desastres de
        infraestructura.
\end{itemize}

\section{Estructura del Equipo de Trabajo y Gobernanza}

LafroX fundamenta su ventaja competitiva en no tercerizar jamás la ingeniería de sus
soluciones. El proyecto está liderado por un grupo multidisciplinario de ocho
Ingenieros Civiles Informáticos con certificaciones internacionales de clase mundial:

\begin{table}[htbp]
\centering
\small
\caption{Roles estratégicos y credenciales del equipo de trabajo}
\label{tab:equipo}
\begin{tabularx}{\textwidth}{@{}L{4.6cm}L{5.4cm}X@{}}
\toprule
\textbf{Rol estratégico} & \textbf{Ingenieros asignados} & \textbf{Función principal} \\
\midrule
Gobernanza y jefatura de proyecto &
Alex Aravena (PM) \newline PMP (PMI), CSM &
Cumplimiento del cronograma maestro, mitigación de la matriz de riesgos, gestión del
presupuesto y comunicación con el Comité Ejecutivo del Mandante. \\
\addlinespace
Arquitectura de soluciones (Cloud \& Edge) &
Tomás Pérez, Álvaro Catalán \newline AWS Certified SA -- Professional &
Diseño conceptual y detallado de la topología lógica-física distribuida, integración
con bases de datos y ERP legados, gobierno de datos y resiliencia en escenarios sin red. \\
\addlinespace
Ingeniería backend e integración &
Leandro Chamorro, Bastián Trejo \newline CKAD &
Desarrollo del catálogo de microservicios, parametrización de motores de reglas de
negocio y optimización de bases de datos relacionales y NoSQL. \\
\addlinespace
Desarrollo frontend y movilidad (offline-first) &
Maximiliano Miño, Guillermo Castillo \newline Mobile UI/UX Specialist &
Diseño e implementación de interfaces de usuario y lógica y bases de datos locales
móviles encriptadas para personal de terreno y conductores. \\
\addlinespace
Ciberseguridad, QA e infraestructura &
Patricio Henríquez \newline CISSP, CEH &
Auditorías de seguridad y penetración, implementación de canalizaciones CI/CD,
simulación de pérdida de red y cumplimiento de protección de datos personales. \\
\bottomrule
\end{tabularx}
\end{table}
